%%%%%%%%%%%%%%%%%%%%%%%%%%%%%%%%%%%%%%%%%%%%%%%%%%%%%%%%%%%%%%%%%%%%%%%%%%%%%%%%%%%%%%
\section{The \primitive primitive}
\label{sec:bbca}
%%%%%%%%%%%%%%%%%%%%%%%%%%%%%%%%%%%%%%%%%%%%%%%%%%%%%%%%%%%%%%%%%%%%%%%%%%%%%%%%%%%%%%

A core building block of \system is \primitive, 
a variant of Byzantine consistent broadcast~\cite{CachinGR11}
with \emph{desired} slot numbers and an external validity predicate $\mathcal{P}$.
The key property we require from \primitive is that slot numbers of delivered blocks are unique across \primitive instances.
Notice that, uniqueness of slot numbers could be expressed as an external validity predicate; 
we make it an explicit part of the primitive specification in order to express properties about the log which is constructed via these slots.

A \primitive instance for some block $b$ is invoked with an
$\bbcacast(id, sn, b, \mathcal{P})$ event, where $id\in\mathbb{N}$ is the instance label
and $sn\in\mathbb{N}$ is the desired slot number.
We say in short that the sender \primitivesmall-bcasts $(id, sn, b)$.
Two distinct instances invoked with the same desired slot number $sn$ are \emph{conflicting}, 
and at most one can succeed in delivering a block with $sn$.

In order to handle faults that potentially leave slots unfilled,
\primitive allows validators to trigger finalization of slot numbers via a \finalize mechanism that triggers an intermediate \bbcaadopt event that has the following guarantee:
If any validator \bbcadeliver's a block $b$ with slot $sn$, 
then $f+1$ correct validators would incur \bbcaadopt of $b$ upon invoking \finalize.
Adopting is supported through a $\bbcaabort(sn)$ interface.
To use the \finalize mechanism, when $\bbcaabort(sn)$ is invoked by a validator,
\primitive stops participating in broadcasts for slot $sn$, and
asynchronously responds to the validator with an adopt event
$\bbcaannounce(id, sn, b, \pi, \mathcal{Q})$ with a (unique) block $b$ (possibly $\bot$). 
$\pi$ is a certificate of uniqueness, and $\mathcal{Q}$ is a predicate indicating whether the block indicated by $\pi$ should be finalized.
We introduce this API event call to allow any \system slot to become finalized within a finite time, potentially with a $\bot$ value.

In order to salvage a \primitive that cannot obtain its desired slot number and uphold broadcast Validity, 
we allow blocks to be delivered without slots, indicated with a $\emptyslot$ slot value.
Delivery events are denoted $\bbcadeliver(id, sn, b)$, where $sn\in\mathbb{N}\cup\{\emptyslot\}$.
We say that the validator \primitivesmall-delivers $(id, sn, b)$.

\Cref{fig:bbca} shows an example with multiple \primitive invocations.

\begin{figure}[h]
  \centering
  \includegraphics[width=\columnwidth]{figures/BBCA}
  \caption{\primitive invocations. Block $c$ is delivered without a slot number, as at most one block can be delivered with slot number $1$.}
  \label{fig:bbca}
\end{figure}


We define \primitive with the following properties:

\begin{primitivelist}
	\item \textit{Integrity}: If a correct validator \primitivesmall-delivers $(id, sn, b)$ and the sender $p$ is correct, then $p$ had previously invoked \primitivesmall-bcast for $(id, sn, b)$.
	\item \textit{No-duplication}: Every correct validator \primitivesmall-delivers at most once per instance.
    \item \textit{Consistency}: If a correct validator $p$ \primitivesmall-delivers $(\bm{id}, sn, b)$ and a correct $q$ validator \primitivesmall-delivers $(\bm{id}, sn', b')$  for the same instance $id$, then $b=b'$ and $sn=sn'$ .
    \item \textit{Unique sequencing}: If a correct validator $p$ \primitivesmall-delivers $(id, \bm{sn}, b)$ and a correct validator $q$  \primitivesmall-delivers $(id', \bm{sn}, b')$ for the same non-nil slot number $sn \in \mathbb{N}$, then $id=id'$ (and by Consistency also $b=b'$).
	\item \textit{Validity}: If a correct validator $p$ \primitivesmall-bcasts $(id, sn, b)$ then eventually every correct validator \primitivesmall-delivers $(id, sn', b)$, such that either $sn'=sn \in\mathbb{N}$ or $sn'=\emptyslot$.
	\item \textit{Weak Sequencing Validity}: After GST, if a correct validator $p$ \primitivesmall-bcasts $(id, sn, b)$ with $sn\in\mathbb{N}$, no other validator \primitivesmall-bcasts $(id',sn, b')$,
	and no correct validator incurs $\bbcaabort(sn)$, 
	then every correct validator delivers $(id, sn, b)$.
	\item \textit{External Validity:} If some correct validator \primitivesmall-delivers $(id, sn, b)$ then $\mathcal{P}(b)$ is \true.
	\item \textit{Finalization:} If some correct validator $p$ incurs\\ $\bbcaabort(sn)$ with $sn \in\mathbb{N}$, then $p$ incurs $\bbcaannounce$.
        Moreover, if $2f+1$ validators with indexes in set $I$ incur \\
        $\bbcaannounce(id_i, sn, b_i, \pi_i, \mathcal{Q})$ such that $\mathcal{Q}(\bigcup_{i\in I}\{\pi_i\})=\false$, then no correct node \primitivesmall-delivers $(id, sn, b)$ for any $id$ and block $b$.
\end{primitivelist}


In \Cref{alg:bbca-pt1} we describe a protocol that implements \primitive.
The protocol makes use of a point-to-point communication primitive, $\plsend(p,b)$ and $\pldeliver(p,b)$,
corresponding to a process sending a block $b$ to process $p$ and to a process receiving a block $b$ from process $p$ respectively.
Likewise, it makes use of a validated Byzantine reliable broadcast primitive,
$\brbcast(b,\mathcal{Q})$ and $\brbdeliver(p,b,\mathcal{Q})$, where
a block $b$ with a dedicated sender $p$ is broadcast with an external validity predicate $\mathcal{Q}$ (see~\cite{CachinKPS01} for definition).
We assume that all events in \Cref{alg:bbca-pt1} are handled atomically.

The core of the \primitive protocol borrows from Bracha's known scheme~\cite{BrachaT85}.
However, the key difference in our protocol is that if validators cannot deliver a block $b$ with its desired slot number $sn$ within a time period, they must 
allow the sender to broadcast $b$ without a slot number.
To this end, they ``convince'' an honest sender that it is safe to rebroadcast $b$ with $sn=\emptyslot$.
Conversely,
if it is possible that some correct validator has already delivered $b$ with $sn\neq\emptyslot$, a correct sender will convince all correct validators to deliver the same way.

\paragraph{The protocol.}
A dedicated \primitive sender $s$ for instance $id$ sends to all validators an $\langle\initiate, id, sn, b, \sigma_s\rangle$ message.
Upon receiving an \initiate message for slot number $sn$ and block $b$ and upon the validity predicate $\mathcal{P}$ for $b$ becomes true, a correct validator $p$ verifies that it is sent by the dedicated instance sender, and that the signature on it is valid.
Moreover, $p$ verifies that it has not echoed another message for either the same instance $id$ or $sn$.
Then $p$ sends a message $\langle\echo, id, sn, b, \sigma_p\rangle$ to all validators and starts a timer $T_{\mathit{id}}$ for instance $id$.

Upon receiving $2f+1$ matching valid \echo messages from different validators, $q$ sends a $\langle\ready, id, sn, b\rangle$ message to all validators.
After receiving $2f+1$ \ready messages,
$q$ \primitivesmall-delivers $id,sn,b$ and cancels the timer.

\Cref{fig:bbca-execution}(a) illustrates the common-case execution of our \primitive implementation.

\paragraph{Finalizing a slot number.}
Validators must eventually finalize each slot, potentially with an empty block, 
in order to allow higher slots to commit. 
To this end, 
when the timer of a slot number $sn$ goes off at a correct validator $p$ it invokes $\bbcaabort(sn)$.
After $\bbcaabort(sn)$ is invoked at $p$, $p$ stops accepting and sending \initiate, \echo and \ready messages for $sn$,
and triggers a $\bbcaannounce(sn, b, \pi, \mathcal{Q})$ event, where $\pi=C_R$,
a set of the $2f+1$ \echo messages $p$ has received in case it has sent a \ready message for slot number $sn$,
or an empty set otherwise.
Notice that there exists at most one instance for which $C_R$ could be non-empty.
Contrary to yield, finalizing a slot number requires a consensus decision, which is described in the next section. 

\paragraph{Yielding the slot number of a \primitive instance.}
If the timer of an instance $id$ goes off, $q$ must ``yield'' the slot number of the instance.
To this effect $q$ stops accepting or sending \initiate/\echo/\ready messages for this instance and sends $\langle\abort, id, sn, b, C_R, \sigma_p \rangle$ to all validators.
$C_R$ is a set of the $2f+1$ \echo messages $q$ has received that triggered \ready,
$C_R = \emptyset$ if none.
We refer to $C_R$ as \emph{ready certificate}.
Moreover, upon receiving $f+1$ matching \abort messages from different validators, if $q$ has not already yielded the slot number of the instance, it does so, too.

When a sender of instance $id$ receives $2f+1$ \abort messages for $id$, it must decide how to help all correct validators \bbcadeliver the corresponding block $b$.
If there exists at least one $C_R \neq \emptyset$ with some $sn$ among the \abort messages, it is safe to assign slot number $sn$ to $b$. By quorum intersection, no validator could have received enough \echo messages for a different block $b'$ carrying the same $sn$.
Therefore, the sender invokes $\brbcast((id,sn,\mathcal{C},b),\mathcal{Q})$, where $\mathcal{C}$ is the set of the $2f+1$ ready certificates and the external validity predicate $\mathcal{Q}$ verifies that indeed the set $C$ justifies assigning $sn$ to $b$.

If, on the other hand, all $C_R$ sets are empty, no validator can \primitivesmall-deliver $b$, as not enough validators will ever send a \ready message.
Therefore, it is safe for the sender to rebroadcast $b$ with no slot number.  
It invokes a $\brbcast((id,\emptyslot,\mathcal{C},b),\mathcal{Q})$, where the external validity predicate $\mathcal{Q}$ checks that indeed the set $C$ justifies assigning the special slot number $\emptyslot$ to $b$.

\Cref{fig:bbca-execution}(b-c) illustrate potential yielding scenarios.

\begin{figure}[h]
  \centering
  \begin{subfigure}{\linewidth}
	  \includegraphics[width=\columnwidth]{figures/BBCA_common_case}  
 	 \caption{Common case execution.}
 	 \label{fig:bbca-common}
  \end{subfigure}
  \begin{subfigure}{\linewidth}
	  \includegraphics[width=\columnwidth]{figures/BBCA_fault_1}  
 	 \caption{Yield where no node sent a \ready message: validator $0$ must invoke \brb with $\emptyslot$.}
 	 \label{fig:bbca-fault-1}
  \end{subfigure}
   \begin{subfigure}{\linewidth}
	  \includegraphics[width=\columnwidth]{figures/BBCA_fault_2}  
 	 \caption{Yield where validator $2$ sent \ready with $sn$: validator $0$ must invoke \brb with $sn$.}
 	 \label{fig:bbca-fault-2}
  \end{subfigure}
  \caption{\primitive executions with and without a timeout for $4$ validators. The dashed line represents the timer.}
  \label{fig:bbca-execution}
\end{figure}

\paragraph{Validity.}
\primitive is weaker than consensus, since
we only require to guarantee Validity for a correct sender.
Moreover, even for a single, correct sender, \primitive requires a checkpointing mechanism to satisfy the Validity property for all correct validators.
In particular, note that a correct validator either delivers a block for some instance, or times out.
It is, however, possible that fewer than $f+1$ validators time out for some instance and, therefore, the sender never triggers \brb to help them \primitivesmall-deliver some block for this instance.
To ensure that all correct validators \primitivesmall-deliver a block for any instance with a correct sender, a checkpointing mechanism should be in place.
A correct validator $p$ sends a $\langle\checkpoint, id, sn, b, \sigma_p\rangle$ message upon \primitivesmall-deliver.
A validator may \primitivesmall-deliver $(id, sn, b)$ upon receiving $f+1$ \textsc{checkpoint} messages, as at least one of them is from a correct validator.
Importantly, this mechanism does not increase latency in the common case and, moreover, validators can batch \textsc{checkpoint} messages.


\begin{algorithm}
  \caption{\primitive implementation.}
  \label{alg:bbca-pt1}
  \scriptsize

  \label{alg:sb-impl}
  \begin{algorithmic}[1]
    \Implements
      \State \primitive \textsc{\primitivesmall} with external validity predicates $\mathcal{P}$ and $\mathcal{Q}$.
    \EndImplements
    \Uses
      \State Byzantine reliable broadcast, \textsc{\brbsmall}
      \State Point-to-point links, \textsc{pl}
    \EndUses
    
    \Parameters
      \State $\numprocesses$
      \Comment{Number of nodes}
      \State $\numfaults$
      \Comment{Maximum number of faulty nodes}
      \State $N$
      \Comment{Set of all $\numprocesses$ nodes}
      \State $\Delta$ 
      \Comment{Progress timeout period}
      \State $T\gets\{\}$
      \Comment{Map of timers per instance label}
      \State \textit{sentecho}$\gets\{\}$
      \Comment{Map of boolean values per instance label}
      \State \textit{sentready}$\gets\{\}$
      \Comment{Map of boolean values per instance label}
      \State \textit{delivered}$\gets\{\}$
      \Comment{Map of boolean values per instance label} 
      \State \textit{yielded}$\gets\{\}$
      \Comment{Map of boolean values per instance label} 
      \State \textit{echos}$\gets\{\}$
      \Comment{Map of maps of $(sn,b)$ tuples per sender per instance label}
      \State \textit{readys}$\gets\{\}$
      \Comment{Map of maps of $(sn,b)$ tuples per sender per instance label}
      \State \textit{yields}$\gets\{\}$
      \Comment{Map of maps of $(sn,b)$ tuples per sender per instance label}
      \State \textit{checkpoints}$\gets\{\}$
      \Comment{Map of maps of $(sn,b)$ tuples per sender per instance label}
      \State \textit{readyCertPerId}$\gets\{\}$
      \Comment{Map ready certificates per sender per instance label}
      \State \textit{readyCertPerSn}$\gets\{\}$
      \Comment{Map ready certificates per sender per sequence number}
      \State \textit{seqnos}$\gets\{\}$
      \Comment{Map of boolean values per sequence number}
      \State \textit{finlalized}$\gets\{\}$
      \Comment{Map of boolean values per sequence number}
    \EndParameters

	\UponEvent{\bbcacast(id, sn, b, \mathcal{P})}
	\Comment{Triggered only by the dedicated sender $s$ of instance with label $id$}
		\If{\textit{finalized}$[sn]=\false$}
		\State $\sigma_s\gets\operatorname{sign}(s,(id, sn, b))$
		\ForAll{q}{N}
			\State\textbf{trigger} $\plsend(q, \langle \initiate, id, sn, b, \sigma_s \rangle)$
		\EndForAll
		\EndIf
	\EndUponEvent

	\UponEvent{\pldeliver(q, \langle \initiate, id, sn, b, \sigma_q\rangle)  \textbf{ and } \mathcal{P}(b)=\true}
	\Comment{Executed by node $p$}
		\If{$q$ is the sender of instance with label $id$}
			\State $T[id].\operatorname{settimer}(\Delta)$
        		\Comment{Set a progress timer}
			\If{\textit{sentecho}$[id]=\false$ \textbf{and} \textit{yielded}$[id]=\false$ \textbf{and} \textit{finalized}$[sn]=\false$ \textbf{and} \textit{seqnos}$[sn]=\false$\textbf{and} $\operatorname{valid}(q,(id, sn, b), \sigma_q)$}
				\label{ln:integrity}
				\State\textit{sentecho}$[id]\gets\true$
				\State\textit{seqnos}$[sn]\gets\true$
				
        		\State $\sigma_p\gets\operatorname{sign}(p,(id, sn, b))$
 				\ForAll{q}{N}
					\State\textbf{trigger} $\plsend(q,\langle \echo, id, sn, b, \sigma_p\rangle)$
				\EndForAll
			\EndIf
		\EndIf
	\EndUponEvent
	
	\UponEvent{\pldeliver(q, \langle \echo, id, sn, b, \sigma_q\rangle)}
	\Comment{Executed by node $p$}
		\If{\textit{yielded}$[id] = \false$  \textbf{and} \textit{finalized}$[sn]=\false$ \textbf{and} $\operatorname{valid}(q,(id, sn, b), \sigma_q)$ \textbf{and} \textit{echos}$[id][q]=\bot$}
			\State\textit{echos}$[id][q]\gets (sn,b) $
		\EndIf
	\EndUponEvent	
	
	\UponExists{(id,sn,b)}{|\{p\in N | \operatorname{\textit{echos}}[id][p]=(sn,b)\}|>2f \operatorname{\textbf{and}}  \operatorname{\textit{sentready}[id]=\false }}
	\Comment{Executed by node $p$}
		\State\textit{sentready}$[id]\gets\true$
		\State \textit{readyCertPerId}$\gets\operatorname{\textit{echos}}[id]$
		\State \textit{readyCertPerSn}$\gets\operatorname{\textit{echos}}[id]$
		\State $\sigma_p\gets\operatorname{sign}(p,(id, sn, b))$
		\ForAll{q}{N}
			\State\textbf{trigger} $\plsend(q,\langle \ready, id, sn, b, \sigma_p\rangle)$
		\EndForAll
	\EndUponExists
	
	\UponEvent{\pldeliver(q, \langle \ready, id, sn, b, \sigma_q\rangle)}
	\Comment{Executed by node $p$}
		\If{\textit{yielded}$[id] = \false$  \textbf{and} \textit{finalized}$[sn]=\false$ \textbf{and} $\operatorname{valid}(q,(id, sn, b), \sigma_q)$ \textbf{and} \textit{readys}$[id][q]=\bot$}
			\State\textit{ready}$[id][q]\gets (sn,b) $
		\EndIf
	\EndUponEvent
	
	\UponExists{(id,sn,b)}{|\{p\in N | \operatorname{\textit{readys}}[id][p]=(sn,b)\}|>2f \operatorname{\textbf{and}}  \operatorname{\textit{delivered}[id]=\false }}
	\label{ln:no-duplication-1}
	\Comment{Executed by node $p$}
		\State\textit{delivered}$[id]\gets\true$
		\State $T[id].\operatorname{cancel}()$
		\Comment{Cancel progress timer with with timeout duration $\Delta$}
		\State\textbf{trigger} $\bbcadeliver(p, id, sn, b)$
		\State $\sigma_p\gets\operatorname{sign}(p,(id, sn, b))$
		\ForAll{q}{N}
			\State\textbf{trigger} $\plsend(q,\langle \checkpoint, id, sn, b, \sigma_p\rangle)$ 
		\EndForAll
	\EndUponExists
	
		\algstore{bbca}
  \end{algorithmic}
\end{algorithm}

\begin{algorithm}
    \scriptsize
    \ContinuedFloat
    \begin{algorithmic}[1]
    \algrestore{bbca}
    
    \UponEvent{T[id].\operatorname{timeout}}
	\Comment{Executed by node $p$}
		\State\textit{yielded}$[id]\gets\true$
		\State $C_R\gets\operatorname{\textit{readyCertPerId}}[id]$
		\State $\sigma_p\gets\operatorname{sign}(p,(id, sn, b, C_R))$
		\ForAll{q}{N}
			\State\textbf{trigger} $\plsend(q,\langle \abort, id, sn, b, C_R,\sigma_p\rangle)$ 
		\EndForAll
	\EndUponEvent
	
	\UponEvent{\bbcaabort(sn)}
	\Comment{Executed by node $p$}
		\State\textit{finalized}$[sn]\gets\true$
		\State $C_R\gets\operatorname{\textit{readyCertPerSn}}[sn]$
		\State\textbf{trigger} $\bbcaannounce(p, sn, b, C_R)$ 
		\label{ln:announcement}
	\EndUponEvent
	
	\UponEvent{\pldeliver(q, \langle \abort, id, sn, b, C_R, \sigma_q\rangle)}
	\Comment{Executed by node $p$}
		\If{$\operatorname{valid}(q,(id, sn, b, C_R), \sigma_q)$ \textbf{and} \textit{yields}$[id][q]=\bot$}
			\State\textit{yields}$[id][q]\gets sn, b, C_R $
		\EndIf
	\EndUponEvent
	
	\UponExists{(id,sn,b)}{|\{p\in N | \operatorname{\textit{yields}}[id][p]=(sn,b,\_)\}|>2f \operatorname{\textbf{and}}  \operatorname{\textit{delivered}[id]=\false }}
	\Comment{Executed by sender $s$}
	\State $\mathcal{C}\gets$ \textit{yields}$[sn]$
		\State \textbf{if exists} $(sn,b,C_R)\in\mathcal{C}$  \textbf{such that} $C_R\neq\emptyset$ \textbf{then}
		\State \hskip \algorithmicindent \textbf{trigger} $\brbcast((id, sn, b,\mathcal{C}),\mathcal{Q})$
				\State{else}
		\State \hskip \algorithmicindent \textbf{trigger} $\brbcast((id, \emptyslot, b, \mathcal{C}), \mathcal{Q})$
		\State{end if}
	\EndUponExists
	
	\UponExists{(id,sn,b)}{|\{p\in N | \operatorname{\textit{yields}}[id][p]=(sn,b,\_)\}|>f \operatorname{\textbf{and}}  \operatorname{\textit{yielded}[id]=\false }}
				\State\textit{yielded}$[id]\gets\true$
		\State $C_R\gets\operatorname{\textit{readyCertPerId}}[id]$
		\State $\sigma_p\gets\operatorname{sign}(p,(id, sn, b, C_R))$
		\ForAll{q}{N}
			\State\textbf{trigger} $\plsend(q,\langle \abort, id, sn, b, C_R, \sigma_p\rangle)$ 
		\EndForAll
	\EndUponExists

	
	\UponEvent{\brbdeliver(p, (id, sn, b), \mathcal{Q})}
		\If{\textit{delivered}[id]=\false }
		 \label{ln:no-duplication-2}
			\State\textit{delivered}$[id]\gets\true$
			\State\textbf{trigger} $\bbcadeliver(p, id, sn, b)$
			\State $\sigma_p\gets\operatorname{sign}(p,(id, sn, b))$
			\ForAll{q}{N}
				\State\textbf{trigger} $\plsend(q,\langle \checkpoint, id, sn, b, \sigma_p\rangle)$ 
			\EndForAll
		\EndIf
	\EndUponEvent
	
	\UponEvent{\pldeliver(q, \langle \checkpoint, id, sn, b, \sigma_q\rangle)}
	\Comment{Executed by node $p$}
	\If{$\operatorname{valid}(q,(id, sn, b), \sigma_q)$ \textbf{and} \textit{checkpoints}$[id][q]=\bot$}
		\State\textit{checkpoints}$[id][q]\gets (sn,b)$
	\EndIf
	\EndUponEvent	
	
	\UponExists{(id,sn,b)}{|\{p\in N | \operatorname{\textit{checkpoints}}[id][p]=(sn,b)\}|>f \operatorname{\textbf{and}}  \operatorname{\textit{delivered}[id]=\false }}
		\State\textit{delivered}$[id]\gets\true$
		\State\textbf{trigger} $\bbcadeliver(p, id, sn, b)$
	\EndUponExists
	
	\caption{\primitive implementation (continues).}	
  	\label{alg:bbca-pt2}
   \end{algorithmic}
\end{algorithm}  


